\documentclass{cv}

\usepackage[utf8]{inputenc}
\usepackage[ngerman]{babel}
\usepackage{hyperref}

\hypersetup{
	colorlinks=true,
	urlcolor=myblue,
	linkcolor=myblue,
	citecolor=myblue,
	pdftitle={Dr. André Kaderli -- Lebenslauf (Professur Angewandte Mathematik)},
	pdfauthor={André Kaderli}
}

% Rote Platzhalter \tbd{...} vor dem Versand ausfuellen oder entfernen.

\begin{document}

	\name{Dr. André Kaderli}
	\tagline{Promovierter theoretischer Physiker (summa cum laude) an der Schnittstelle von mathematischer Forschung, praktischer Anwendung und Lehre. Seit über fünf Jahren setze ich mathematische Modellierung, Statistik und Machine Learning in Projekten für öffentliche Verwaltung und Privatwirtschaft um, seit 2026 als Partner mit Führungsverantwortung. Mit über 15 Jahren Lehrerfahrung -- von Hochschulübungen bis zu Weiterbildungskursen für Fachleute -- verbinde ich Theorie und Praxis in einem anwendungsnahen Unterricht.}
	\begin{contactline}
		Obere Bahnhofstrasse 22M, 3700 Spiez\sep
		\href{mailto:andre.kaderli@gmail.com}{andre.kaderli@gmail.com}\sep
		\href{tel:+41793323931}{+41\,79 332 39 31}\sep
		\href{https://www.linkedin.com/in/andre-kaderli/}{linkedin.com/in/andre-kaderli}
	\end{contactline}

	\datedsection{Lehrgebiete}{}
	\kv{Mathematische Grundlagen:}{Analysis, Lineare Algebra, Funktionentheorie, Wahrscheinlichkeitsrechnung und Statistik}
	\kv{Angewandte Mathematik:}{Numerik, Optimierung, mathematische Modellierung und Simulation, stochastische Prozesse, Zeitreihenanalyse, Signalverarbeitung}
	\kv{Data Science \& KI:}{Statistisches Lernen und Machine Learning, Datenanalyse und Visualisierung mit Python und R, Grundlagen von KI- und LLM-Anwendungen}

	\datedsection{Lehre und Betreuung}{}

	\datedsubsection{Data Academy, Novalytica}{Seit \tbd{Monat} 2022}
	\desc{Konzeption und Durchführung von Kursen zu Data Literacy, Programmierung (Python, R), Statistik, Datenanalyse und Visualisierung für Fachleute aus Verwaltung und Wirtschaft; Kundenschulungen zu eingeführten Datenplattformen (Microsoft Fabric, Power BI, Databricks); rund \tbd{X} Kurse pro Jahr}

	\datedsubsection{Vorlesungsunterstützung, HU Berlin}{2018 -- 2021}
	\desc{Funktionentheorie und Quantenfeldtheorie (\tbd{X} Semester); Betreuung von Studierenden und wissenschaftlichen Arbeiten in Mathematik und Physik (\tbd{X} Bachelor-/Masterarbeiten)}

	\datedsubsection{Teaching Assistant, ETH Zürich}{2014 -- 2016}
	\desc{Übungsgruppen und Betreuung der Studierenden in Linearer Algebra am Departement Mathematik (\tbd{X} Semester, ca.~\tbd{X} Studierende pro Gruppe)}

	\datedsubsection{Privater Nachhilfelehrer}{2009 -- 2018}
	\desc{Mathematik, Physik und Informatik für Lernende verschiedener Stufen}

	\newpage
	\datedsection{Berufspraxis}{}

	\datedsubsection{Novalytica AG, Bern}{}

	\datedsubsubsection{Principal Data \& AI Engineer und Partner}{Seit Jan. 2026}
	\begin{itemize}
		\item Fachliche Leitung von Daten- und KI-Initiativen in öffentlicher Verwaltung und Privatwirtschaft; Führung und Mentoring interdisziplinärer Teams (\tbd{X} Personen)
		\item Konzeption und Aufbau moderner Daten- und Analytics-Plattformen und Data-Governance-Frameworks; Beratung von Entscheidungsträgern
	\end{itemize}

	\datedsubsubsection{Senior Data Scientist}{Jan. 2023 -- Dez. 2025}
	\begin{itemize}
		\item Mathematische Modellierung komplexer Systeme und statistische Analysen grosser, heterogener Datensätze als Grundlage für Entscheidungen in Verwaltung und Wirtschaft
		\item Leitung von End-to-End-Projekten (\tbd{X} Projekte): von der Problemformulierung über Prognose- und ML-Modelle bis zum produktiven Betrieb; Verantwortung für Entwicklungsstandards (CI/CD, Testing, Monitoring)
	\end{itemize}

	\datedsubsubsection{Data Scientist}{Mai 2022 -- Dez. 2022}
	\begin{itemize}
		\item Entwicklung von ML-Modellen und statistischen Auswertungen; Data-Analytics- und BI-Projekte (Power BI, Azure, Databricks)
	\end{itemize}

	\datedsubsection{Eraneos AG, Bern}{}
	\datedsubsubsection{Consultant}{Sep. 2021 -- Apr. 2022}
	\begin{itemize}
		\item IT- und Datenprojekte im öffentlichen Sektor mit Fokus auf Business-Analyse und quantitative Datenauswertungen
	\end{itemize}

	\datedsection{Forschung}{}

	\datedsubsection{Max-Planck-Institut für Gravitationsphysik, Potsdam}{}
	\datedsubsubsection{Wissenschaftlicher Mitarbeiter (Doktorand)}{Sep. 2018 -- Aug. 2021}
	\begin{itemize}
		\item Forschung zur Quantengravitation an der Schnittstelle von Stringtheorie und Zahlentheorie (u.\,a. elliptische Multi-Polylogarithmen)
		\item Analytische und numerische Methoden, rekursive Algorithmen und Simulationsverfahren zur Lösung komplexer Differentialgleichungen
		\item Konferenzbeiträge und Vorträge: \tbd{Anzahl, ggf. wichtigste Konferenzen}
	\end{itemize}

	\datedsubsection{Ausgewählte Publikationen}{}
	\desc{\tbd{Autoren, «Titel», Zeitschrift Band (Jahr), Seiten}}
	\desc{\tbd{Autoren, «Titel», Zeitschrift Band (Jahr), Seiten}}
	\desc{\tbd{Autoren, «Titel», Zeitschrift Band (Jahr), Seiten}}
	\desc{Vollständige Liste: \href{https://inspirehep.net/authors/1809024}{inspirehep.net/authors/1809024}}

	\newpage
	\datedsection{Ausbildung}{}

	\datedsubsection{PhD in Physik, Humboldt-Universität zu Berlin}{Sep. 2018 -- Aug. 2021}
	\desc{Theoretische Physik (summa cum laude); Dissertation: «\tbd{Titel}», Betreuung: Prof.~\tbd{Name}}

	\datedsubsection{MSc in Physik, ETH Zürich}{Sep. 2016 -- Aug. 2018}
	\desc{Theoretische und Mathematische Physik (Abschlussnote: 6.0/6.0); Fokus auf Allgemeine Relativitätstheorie, Quantenfeldtheorie und Statistische Physik}

	\datedsubsection{BSc in Physik, ETH Zürich}{Sep. 2013 -- Aug. 2016}
	\desc{Physik, Mathematik und Informatik (Abschlussnote: 5.7/6.0)}

	\datedsection{Fachliche Kenntnisse}{}
	\kv{Mathematik:}{Lineare Algebra, Analysis, Funktionentheorie, Wahrscheinlichkeitstheorie, Statistik, Algebra, Zahlentheorie, Geometrie, Differentialgeometrie, Topologie, algebraische Topologie}
	\kv{Physik:}{Klassische Mechanik, Elektrodynamik, Thermodynamik und Statistische Physik, Quantenmechanik, Allgemeine Relativitätstheorie, Quantenfeldtheorie, Stringtheorie, Festkörperphysik, Quantum Computing}
	\kv{Programmierung \& Werkzeuge:}{Python, R, C++, MATLAB, Mathematica, SQL, PySpark; Azure, Databricks, Microsoft Fabric, Power BI}

	\datedsection{Weitere Qualifikationen}{}
	\kv{Zertifizierungen:}{Microsoft Azure Data Scientist, Azure Data Engineer, Power BI Data Analyst, Fabric Analytics Engineer und Fabric Data Engineer (je Associate); Databricks Data Engineer Professional und Associate}
	\kv{Projektmanagement:}{Hermes Advanced}

	\datedsection{Weitere Angaben}{}
	\kv{Sprachen:}{Deutsch (Muttersprache), Englisch (verhandlungssicher), Französisch (Grundkenntnisse)}
	\kv{Stiftungen \& Stipendien:}{Alumni der Schweizerischen Studienstiftung, ehem. Träger des Ernst-Göhner-Stipendiums}
	\kv{Freiwilliges Engagement:}{Tourenleiter beim Schweizer Alpen-Club (SAC)}
	\kv{Interessen:}{Bergsteigen, Klettern, Skifahren, Gleitschirmfliegen, Jiu-Jitsu}

\end{document}
